\section{Dante's Contemplation}

\begin{quotex}
What is a worse punishment than always to will what will never be and to be constantly opposing what will always be? … The soul will never have what it wants and will endure forever what it does not want. \flright{\textsc{Bernard of Clairvaux}}

\end{quotex}
Besides \textbf{Saint Augustine}'s \emph{De Quantitate Anima}, \textbf{Dante} relied heavily on two other works. One is \emph{De Contemplatione} (On Contemplation) by \textbf{Richard of Saint Victor} and the other is \emph{De Consideratione} (On Consideration) by \textbf{Saint Bernard of Clairvaux}. Moreover, it was Bernard who guided Dante to Mary and then to the very Vision of God. The Paulist edition of St. Bernard asserts that ``Dante chose Bernard as his spiritual guide", as though it were an arbitrary literary device. But those who know, understand that it was Bernard who chose Dante. Follow the path described by Bernard and then you will see.

\paragraph{Contemplation}
In his letter to Con Grande\footnote{\url{https://faculty.georgetown.edu/jod/cangrande.english.html}}, explaining the proper interpretation of the \emph{Divine Comedy}, Dante asserted that in contemplation, the intellectual soul is raised to the extent that it transcends human kind. Richard of St. Victor is claimed as a witness to this state, along with Bernard and Augustine. Richard mentions the three ways of thinking: Cogitation, Meditation, and Contemplation. He describes them like this:

\textbf{Cogitation} rambles indifferently here and there in a haphazard way, without regard for the outcome of its rambling. Cogitation creeps. Cogitation involves no labour and brings no fruit. Cogitation issues from imagination. In cogitation is prone to aimless wandering.

\textbf{Meditation} strives with great assiduity, often over a harsh and difficult ground, to reach the end of the road. Meditation walks and runs. In meditation, there is labour and fruit. Meditation issues from reason. In meditation, we scrutinize. It is an eager and persistent effort of the mind in searching and finding out; or a steadfast and careful speculation of the mind passionately fixed on the search of truth.

\textbf{Contemplation} is, with admirable agility, borne up in its free flight, wheresoever inspiration carries it. Contemplation arises from intelligence. Contemplation is an unlimited clarity of the mind, suspended in admiration at the sight of wisdom. Thus, it is an unrestricted, all-embracing penetration of the mind into those things, which are to be comprehended.

\begin{quotex}
These are the three sources affiliated with those ways. \textbf{Intelligence} holds the highest place, \textbf{imagination} the lowest and \textbf{reason} is in the middle. Every notion, which is subject to the inferior mode of thinking, must also be subject to the superior one. It follows that ideas, grasped by imagination, and many others above them, are grasped by reason as well. Similarly, those that are apprehended by imagination or by reason, as well as those that reason and imagination cannot apprehend, submit to intelligence. See then how widely the radius of contemplation expands, how it embraces and lights up everything. 

\end{quotex}
Back to the epigraph: \textbf{Imagination} fantasizes a future that will never be. Even worse, it fantasizes a past that could have been, but never was. The solution is to evaporate such fantasies under the bright light of consciousness.

\textbf{Reason}, on the other hand, when it is dissatisfied with what is, tries to think of an alternative. The dream of reason is to resolve all the questions of life through philosophy, science, drugs, psychological therapies, five-year programs, ``new" deals, and so on. Although it may ameliorate some harsh conditions of life, the fundamental human problems of ignorance, malice, concupiscence, and death remain intractable.

The second punishment, which arises from reason, is to rue and lament what actually is. The antidote is abandonment to divine providence, or nihilism.

\textbf{Intelligence} is a higher order thinking that transcends the human state. The mind becomes clear, the soul is set at ease, and the intellect has a direct grasp of reality.

\paragraph{Consideration}
\begin{quotex}
The greatest of all is he who spurns the use of things which the senses can perceive and goes up not by steps but in great leaps beyond our imagining: he has learned to fly to the heights in contemplation at times. \flright{\textsc{Bernard of Clairvaux}}

\end{quotex}
Prior to the greatest consideration, which is the purest, there are two others, one powerful and the other freer. Bernard gives names to the three types of consideration:

\begin{enumerate}
\item \textbf{Practical}: It makes use of the senses and the things the senses perceive in an orderly and coordinated way, so as to please God. 
\item \textbf{Scientific}: It wisely and carefully searches into and weighs the signs of God's work in the world. 
\item \textbf{Speculative}: It retires into itself and, as far as God helps it, frees itself from human affairs for the contemplation of God. 
\end{enumerate}
Although each of the forms is suited to those in his particular station in life, the third is the fruit of the other two, which otherwise meander aimlessly without the third in view. There are three ways to pursue knowledge, from the lowest to the highest:

\begin{enumerate}
\item \textbf{Opinion}. Opinion claims no certainty because it tries to discover what is like truth, rather than to grasp it directly. It considers something to be true that is not known to be false. Eight centuries before Karl Popper, the Scholastics understood that opinion could not be proven true, but only falsified. 
\item \textbf{Faith}. Faith is a certain voluntary and confident foretaste of truth not yet apparent. However, truth is hidden and obscure. 
\item \textbf{Gnosis}. Gnosis is a clear and certain grasp of something unseen. Gnosis not only grasps the truth, but knows that it is the truth. 
\end{enumerate}
Opinion is necessarily provisional, but to believe something false is ignorance. Faith has no doubt, otherwise it is opinion. When it is true, faith is still a mystery, since it lacks gnosis. Then you know something, you seek no further; otherwise, it is not true gnosis\footnote{Bernard goes on to suggest some topics worthy of contemplation. The first of those is the hierarchy of angels, which has special relevance to the \emph{Divine Comedy}. That will be addressed in a post in the near future.}.

\begin{quotex}
Our happiness will be complete when what is already certain to us will be as plain as it is certain. 

\end{quotex}

\flrightit{Posted on 2020-08-23 by Cologero }

\begin{center}* * *\end{center}

\begin{footnotesize}\begin{sffamily}



\texttt{Lee on 2020-08-24 at 02:19 said: }

No, I would say the concept of gnosis is ultimately outdated. First of all that is a semantical maneuver that makes the concept appear more appealing than it really is. Secondly, advances such as pragmatism show that gnosis ultimately leads to error. For instance, as pragmatism states, there are contradictory claims within gnosis itself. For instance the yin and the yang comes in contradiction with the angelic hierarchy, which comes in contradiction with the law of attraction. Ultimately transcendent claims should be attempted to be rationalized as much as possible. This allowsiallows for the most room for original work. how can the yin and the yang possibly be true given the contradiction with the sciences. This is one of the strengths of modern idealism. Once you allow for gnosis you cut off the Publics can rationality which can be dangerous. The modified argument is that gnosis can never be used never to contradict another theory. You need the give up the critique of posterior reasoning as given by Kant.


\hfill

\texttt{Cologero on 2020-08-24 at 05:15 said: }

It's outdated, Lee? So when will your opinion be outdated? Can you predict the most compelling argument in the year 2525, if man is still alive?


\hfill

\texttt{Dennis on 2020-08-24 at 10:14 said: }

``advances such as pragmatism…"

Thanks for the laugh Lee…needed that on a Monday morning!


\hfill

\texttt{James on 2020-08-25 at 02:02 said: }

The seeming ``contradictions" in gnosis is not at all a bad thing . It is part of the task of attaining gnosis that one is able to experience the synthesis of contradictory things . This is attested to by the highest reality: God + Man . The Source of All Life + Death . An instrument of torture , death , and humiliation + paradise , everlasting life , and glory . Three is One and One is Three . A creature being the Mother of the creator . The universal + the particular . Eternity + Time .

The exercise of being able to understand seeming opposites is , in fact , the initiation from above because it is from above that everything that rises converges (to quote that beautiful Mystic Flannery O'Connor as she herself gets it from Teilhard de Chardin) . To be able to retrace the rainbow back to the pure white light .

Contradictions , at least in the sphere of gnosis and not simply the cult of chaos , is actually a participation in the creative act of the universe and of man . For creation is the synthesis of two seemingly contradictory things: nothingness and being . Chaos and Order . Choice and Divine Will . This is why in order to start on the path to gnosis , one must first find that ``rapport between individual effort and divine reality" as a dear friend once said . It is the ability to calmly , and in silence , attain to the ``marriage of opposites" . As man and woman as seeming opposites unite , so does a new life come forth .

``The ying and yang" being in contradiction to the ``hierarchy of angels" might only seem contradictory when looked upon in purely materialistic terms . One might even say that it's not even fair to say that they are contradictory in exoteric terms ! One need only understand the purpose of each . French is no more valid of a language than English even if English has two words for Beef and Cow . The language may be different , but the meaning might be compatible . They are only inappropriate exoterically in the same way that if one is invited into the house of a host , one should have the courtesy to eat *his* food rather than attempt to override the cook . Nothing is wrong with the food of either (so long as it is conducive to health , of course) , but it is a matter of congeniality that one acknowledges the benefits of the other . 

Transcendent claims can indeed and have been for the longest time rationalized . This has never been in dispute . An ant crawling on the inked manuscript of Shakespeare's sonnets reduces the ink to ``this is a dangerous chemical , i shall not step on this" without any idea that it is some of the greatest lines in English . I do not begrudge the ant for believing that . Her instincts have served her well . She may be the very best ant in the colony .

I congratulate you on taking the time to warn the other people of the dangers of transcendent thinking , Lee . After all , the organism needs feet as much as it needs a heart and and eyes . We should be wounded as creatures if our feet attempted to pump blood or our knees attempted to ``see" . Thus , I encourage you to maintain your vigilant adherence to reduction . It will serve all of us well if that is your true vocation in our society and are incapable of any other type of thought . I look forward to offering you complex logic puzzles to solve so that the rest of us don't have to worry about such tasks .

You believe that originality can be attained by divorcing ourselves from higher thought (or at least convincing ourselves that there is no such thing as `higher' thought) . I do not know what you consider to be the definition of ``originality" .. and I also do question if originality as the end is truly conducive for the humanity of a person , but it does not matter . I'm afraid you cannot save Cologero . He is too far gone . Better to try with others who are more willing to be productive reductionists . Too many of us here are obsessed with , as Oscar Wilde might put it in his Picture of Dorian Grey , ``quite useless" things . But I thank you for your service to our society .


\hfill

\texttt{Matt on 2020-08-25 at 17:51 said: }

``Secondly, advances such as pragmatism show that gnosis ultimately leads to error."

Error flows from knowledge? Is this some attempt at a Zen koan?


\hfill

\texttt{Santiago on 2021-11-16 at 15:37 said: }

One of the most useful summations I have ever read.


\end{sffamily}\end{footnotesize}
